% ===================== CHAPITRE 6 — SPRINT 4 =========================
\chapter{Sprint 4 : Pointage, demandes et notifications}

\section*{Introduction}
Ce sprint regroupe trois mécanismes transverses : le \textbf{pointage} de présence
des employés, les \textbf{workflows de demandes} (création et modification) et le
dispositif de \textbf{notifications} et de \textbf{journalisation} des actions.

\section{Objectifs du Sprint}
L'objectif est d'assurer le \textbf{suivi de présence} des employés, d'\textbf{encadrer}
les opérations sensibles par des circuits de validation et de garantir la
\textbf{traçabilité} des actions par la journalisation et les notifications.

\section{Spécification des besoins}

\subsection{Diagramme de cas d'utilisation du Sprint 4}
La figure~\ref{fig:uc_s4} présente les cas d'utilisation du sprint.

\begin{figure}[htbp]
  \centering
  \includegraphics[width=0.92\textwidth]{uc_sprint4.png}
  \caption{Diagramme de cas d'utilisation du Sprint 4}
  \label{fig:uc_s4}
\end{figure}

\subsection{Sprint Backlog du sprint 4}
Le tableau~\ref{tab:bl_s4} détaille les user stories et les tâches techniques.

\begin{table}[htbp]
\centering
\caption{Sprint Backlog du sprint 4}
\label{tab:bl_s4}
\renewcommand{\arraystretch}{1.3}
\begin{tabularx}{\textwidth}{>{\bfseries}p{1cm}p{5cm}X}
\toprule
\textmd{\textbf{US}} & \textbf{User story} & \textbf{Tâches techniques}\\
\midrule
US9 & Pointer la présence (arrivée / départ). &
Développer l'entité \texttt{Pointage}, le \texttt{PointageService} (EJB) et le
\texttt{PointageServlet} (\texttt{/pointage}) ; déterminer le statut
(\texttt{PRESENT} / \texttt{RETARD} au-delà de 9h) ; page \texttt{pointage.jsp}.\\
\midrule
US10 & Saisir ou corriger un pointage manuel. &
Permettre à l'administrateur de saisir un pointage (\texttt{PRESENT},
\texttt{RETARD}, \texttt{ABSENT}) via \texttt{PointageService.saveManuel}.\\
\midrule
US11 & Soumettre une demande (client, employé, modification). &
Développer les entités \texttt{DemandeClient}, \texttt{DemandeEmploye},
\texttt{DemandeModification}, \texttt{DemandeAgence} et les servlets
\texttt{DemandeAgenceServlet} / \texttt{MesDemandesServlet} ; statut
\texttt{EN\_ATTENTE}.\\
\midrule
US12 & Valider ou refuser une demande. &
Développer le \texttt{GestionModificationsServlet} et le
\texttt{DemandeAgenceAdminServlet} ; faire évoluer le statut vers
\texttt{EFFECTUE} / \texttt{REFUSE} ; pages \texttt{gestion-demandes.jsp},
\texttt{demandes-admin.jsp}.\\
\midrule
US13 & Recevoir des notifications et consulter le journal d'audit. &
Développer le \texttt{NotificationService}, l'entité \texttt{Notification} et le
\texttt{NotificationFilter} ; journaliser les actions via le
\texttt{JournalService} (\texttt{JournalAction}).\\
\bottomrule
\end{tabularx}
\end{table}

\section{Analyse des besoins}

\subsection{Description textuelle des cas d'utilisation}

\cudesc{tab:cu_pointage}{Pointer la présence}
{Permettre à l'employé d'enregistrer son arrivée et son départ, avec
détermination automatique du statut.}
{Utilisateur (employé)}
{L'employé doit être authentifié.}
{Le pointage du jour est enregistré avec son statut (\texttt{PRESENT} ou
\texttt{RETARD}).}
{(1) L'employé clique sur « Pointer » ; (2) le \texttt{PointageService} relève
l'heure courante et recherche le pointage du jour ; (3) s'il n'existe pas, il
détermine le statut (\texttt{RETARD} au-delà de 9h) et l'enregistre ; (4) sinon,
il met à jour l'enregistrement existant.}

\cudesc{tab:cu_demande}{Valider une demande}
{Permettre à l'administrateur d'approuver ou de refuser une demande soumise.}
{Administrateur}
{Une demande au statut \texttt{EN\_ATTENTE} doit exister.}
{Le statut de la demande évolue vers \texttt{EFFECTUE} ou \texttt{REFUSE} et
l'action est journalisée.}
{(1) L'administrateur consulte les demandes en attente ; (2) il sélectionne une
demande ; (3) il l'approuve ou la refuse avec un commentaire ; (4) le service met
à jour le statut, applique la modification le cas échéant et journalise l'action.}

\section{Conception}

\subsection{Diagramme de séquence : « Pointer un employé »}
La figure~\ref{fig:seq_pointage} décrit le scénario de pointage.

\begin{figure}[htbp]
  \centering
  \includegraphics[width=0.95\textwidth]{seq_pointage.png}
  \caption{Diagramme de séquence -- Pointage d'un employé}
  \label{fig:seq_pointage}
\end{figure}

Le \texttt{PointageServlet} transmet l'action au \texttt{PointageService}, qui
recherche le pointage du jour. Le fragment \texttt{alt} distingue l'\texttt{INSERT}
(premier pointage, statut \texttt{PRESENT}/\texttt{RETARD} au-delà de 9h) de
l'\texttt{UPDATE} (pointage déjà existant).

\subsection{Diagramme de séquence : « Valider une demande »}
La figure~\ref{fig:seq_demande} détaille le traitement d'une demande par
l'administrateur.

\begin{figure}[htbp]
  \centering
  \includegraphics[width=0.92\textwidth]{seq_demande.png}
  \caption{Diagramme de séquence -- Validation d'une demande}
  \label{fig:seq_demande}
\end{figure}

Le \texttt{GestionModificationsServlet} transmet la décision au service métier,
qui charge la demande, met à jour son statut, persiste la modification puis
journalise l'action via le \texttt{JournalService}.

\subsection{Diagramme d'états-transitions : cycle de vie d'une demande}
La figure~\ref{fig:state_demande} modélise le cycle de vie d'une demande.

\begin{figure}[htbp]
  \centering
  \includegraphics[width=0.6\textwidth]{state_demande.png}
  \caption{Diagramme d'états-transitions -- Cycle de vie d'une demande}
  \label{fig:state_demande}
\end{figure}

Une demande naît à l'état \texttt{EN\_ATTENTE}, passe en \texttt{EN\_COURS} lors de
sa prise en charge, puis atteint un état final \texttt{EFFECTUE} ou
\texttt{REFUSE}, en cohérence avec la contrainte de la table des demandes.

\subsection{Diagramme d'activité : traitement d'une demande}
La figure~\ref{fig:activity_demande} représente le processus métier de traitement
d'une demande.

\begin{figure}[htbp]
  \centering
  \includegraphics[width=0.6\textwidth]{activity_demande.png}
  \caption{Diagramme d'activité -- Traitement d'une demande}
  \label{fig:activity_demande}
\end{figure}

Le directeur soumet une demande, notifiée à l'administrateur. Le nœud de décision
oriente le flux : en cas de validation, le statut passe à \texttt{EFFECTUE} et la
modification est appliquée ; en cas de refus, le statut passe à \texttt{REFUSE}
avec un commentaire. Les deux branches convergent vers la journalisation et la
notification du demandeur.

\section{Réalisation : interfaces utilisateur}
La figure~\ref{fig:ui_pointage} présente l'interface du module de pointage : les
cartes de synthèse (présents, retards, absents, taux de présence), le pointage du
jour (boutons « Pointer l'arrivée / le départ ») et le formulaire de saisie ou de
correction manuelle.

\begin{figure}[htbp]
  \centering
  \includegraphics[width=\textwidth]{screens/pointage.png}
  \caption{Interface du module de pointage des employés}
  \label{fig:ui_pointage}
\end{figure}

\section*{Conclusion}
Ce sprint a doté l'application du suivi de présence, des circuits de validation et
de la traçabilité des actions. Le sprint suivant met en place la chaîne ETL et le
volet décisionnel.
